Large-language-model multi-agent systems are highly sensitive to how agents communicate, which evidence they may reuse, and when a tentative answer should be forwarded. We introduce PromptNode-Augmented multi-agent graphs, a graph-organization method that treats virtual functional nodes as first-class elements rather than post-hoc answer patches. PromptNodes carry typed identities such as memory, instruction verification, tool-resource checking, and adversarial verification; PromptEdges specify read, write, and forward permissions; and lifecycle actions update, compress, prune, merge, or upgrade graph state during reasoning. Using ARG-Designer as the main baseline and preserving its model setting as closely as possible, the current evidence supports a scoped improvement claim. On the audited AQuA tail split, the PromptNode graph obtains 31/31 against ARG's 28/31. On MMLU STEM transfer, it obtains 95/97 against 88/97. A targeted DefinitionChallenge non-STEM run improves 21/24 to 22/24 with fix/harm/wrong-forward 1/0/0, while a broader mixed non-STEM split remains no-harm rather than positive at 41/48 versus 41/48. These results suggest that typed virtual nodes and permissioned prompt edges can improve robust reasoning when claims are bounded by label audits, cost accounting, and transfer diagnostics.
